\documentclass[11pt,a4paper]{article}

% Русская версия. Собирается XeLaTeX/LuaLaTeX: кириллица через fontspec,
% маркер — отдельным шрифтом, в котором есть и хангыль, и Geometric Shapes.
\usepackage{fontspec}
\setmainfont{Times New Roman}
\setmonofont{Consolas}
% Маркер набирается шрифтом, в котором есть и хангыль, и Geometric Shapes.
% Порядок — по доступности: Windows, затем семейства Noto CJK для Linux.
\newif\ifmarkerfontfound
\markerfontfoundfalse
\IfFontExistsTF{Arial Unicode MS}{%
  \newfontfamily\markerfont{Arial Unicode MS}\markerfontfoundtrue}{}
\ifmarkerfontfound\else
  \IfFontExistsTF{Noto Sans CJK KR}{%
    \newfontfamily\markerfont{Noto Sans CJK KR}\markerfontfoundtrue}{}
\fi
\ifmarkerfontfound\else
  \IfFontExistsTF{Noto Sans CJK JP}{%
    \newfontfamily\markerfont{Noto Sans CJK JP}\markerfontfoundtrue}{}
\fi
\ifmarkerfontfound\else
  \IfFontExistsTF{Malgun Gothic}{%
    \newfontfamily\markerfont{Malgun Gothic}\markerfontfoundtrue}{}
\fi
\ifmarkerfontfound
  \newcommand{\marker}[1]{{\normalfont\markerfont #1}}
\else
  \newcommand{\marker}[1]{\texttt{#1}}
\fi
\usepackage{amsmath,amssymb,amsfonts}
\usepackage{booktabs,array,multirow}
\usepackage{microtype}
\usepackage{graphicx}
\usepackage{cite}
\usepackage{url}
\usepackage{xcolor}
\usepackage[margin=1in]{geometry}
\usepackage{hyperref}
\sloppy

\hypersetup{
    colorlinks=true,
    linkcolor=blue!70!black,
    citecolor=red!70!black,
    urlcolor=blue!70!black
}

\title{\textbf{Закон дисперсии внимания и двухканальное расщепление\\в авторегрессионных трансформерах}}

\author{
    \textbf{Владимир В.~«SkyAngel»~Боровиков}\thanks{\href{https://orcid.org/0009-0005-7773-3581}{ORCID: 0009-0005-7773-3581}} \\
    \textit{Независимый исследователь} \\
    \texttt{skywhiteangels@gmail.com}
}

\date{Октябрь 2026}

\begin{document}

\maketitle

\begin{abstract}
Большие языковые модели хорошо продолжают связный текст, но остаются уязвимы там, где нужно удерживать \textit{дискретные инварианты} — точный целочисленный счёт на длинных последовательностях. На предрегистрированном корпусе из 1\,416 канонических проб ($1\,320$ в основной серии Argon v8--v16 плюс $96$ в репликации v6), снятых при нулевой температуре на современных MoE-трансформерах (при строгой изоляции No-Tools на всех 960 чисто когнитивных пробах счёта в v9 и v11--v16, наряду с протоколами работы с инструментами в v8 и v10), мы исследуем механистический локус ошибки дискретного счёта (перелёт $N \to N+1$ на границе).

Через серию анатомических интервенционных экспериментов мы показываем, что:
\begin{enumerate}
    \item \textbf{Декодер оправдан:} применив алгебраический усилитель ($f(N) = 10N - 6$, v15), чтобы отделить дрейф сканирования внимания ($\Delta = \pm 10$) от глитчей выходного эмбеддинга ($\Delta = +1$), мы получаем строго \textbf{0 / 96} ($0{,}0\%$, двусторонний 95\% интервал Клоппера--Пирсона $[0{,}0\%; 3{,}8\%]$) ошибок декодера на единицу, тогда как на перцептивный дрейф входного сканирования приходится 34{,}4\% (33/96).
    \item \textbf{Гипотеза токенизатора отпадает:} сопоставив побайтово и посимвольно повторяющийся корейский субсловный маркер (\marker{점점점◁}, 12 байт UTF-8) с чистым блоком геометрических символов Unicode (\marker{▤▤▤▤}, 12 байт UTF-8, v16), мы показываем, что замена языка геометрией не устраняет фантомные ошибки ($35{,}4\%$ против $41{,}7\%$, точный тест Фишера $p = 0{,}675$). Это не поддерживает сильную гипотезу BPE-токенизации и согласуется с гипотезой о том, что предел дискретного счёта отражает фазовое ограничение внимания.
    \item \textbf{Двухканальное расщепление внимания:} в состязательных пробах с газлайтингом (v13) модель выполняет указание выдать ложь в видимом содержании ($93{,}8\%$), но спонтанные «утечки истины» напрямую из латентного канала рассуждения (\texttt{thinking}) в терминальный токен происходят в 2/48 пробах. Это показывает, что терминальные логиты могут расходиться с локальным авторегрессионным $n$-граммным контекстом.
\end{enumerate}
Мы формализуем эти эмпирические результаты в \textit{закон дисперсии внимания}: когда монотонные последовательности превышают ёмкость локального окна внимания ($N \ge 8..16$), фазовая интерференция позиционных эмбеддингов размывает границы рецептивных полей и порождает дискретный перелёт. Отсюда следует, что детерминированные агентные архитектуры должны дополнять передачу данных внутри контекста внешней передачей по ссылке и детерминированными рантайм-проверками чётности.
\end{abstract}

\vspace{1em}
\noindent\textbf{Ключевые слова:} интерпретируемость трансформеров, пределы self-attention, двухканальное внимание, индукционные головы, счёт в контексте, безопасность LLM-агентов.

\newpage

\section{Введение}

Авторегрессионные трансформеры произвели революцию в непрерывном моделировании естественного языка, генерации кода и семантическом выводе \cite{vaswani2017attention}. При этом автономные агенты в промышленной эксплуатации регулярно отказывают на простых детерминированных операциях над структурированными строками: подсчёт разделителей, отслеживание состояний автомата, балансировка кавычек, оценка глубины синтаксического дерева.

В промышленном контексте (автоматическая правка кода, разбор бинарных логов, исполнение скриптов) ошибка на единицу ($N \to N+1$ или $N \to N-1$) не является безобидной стилистической погрешностью. Она ведёт к взаимным блокировкам в командных интерпретаторах, сломанным переходам синтаксиса и тихому повреждению данных. В исследованиях LLM сложились две конкурирующие обиходные гипотезы, объясняющие этот сбой счёта:
\begin{itemize}
    \item \textbf{Гипотеза декодера ($H_{\text{DEC}}$):} внутреннее скрытое представление точно кодирует истинный счёт $N$, но финальная линейная проекция (голова эмбеддинга) и температурная дисперсия softmax приводят к случайному перелёту на единицу на выходе.
    \item \textbf{Гипотеза токенизатора ($H_{\text{BPE}}$):} многобайтовые символы и редкие токены подвергаются неравномерному разбиению на субслова при Byte-Pair Encoding (BPE), что вызывает пространственные фазовые сдвиги в ротационных позиционных эмбеддингах (RoPE \cite{su2024roformer}) и вводит модель в заблуждение.
\end{itemize}

\textbf{Экспериментальная серия Argon} была задумана как многоэтапная эмпирическая программа для изоляции, проверки и разрешения этой уязвимости при строгом контроле условий. На 1\,416 канонических пробах мы подвергли современные рантаймы трансформеров строгим интервенционным пробам.

\subsection{Смежные работы: схемы внимания и перцептивные иллюзии}

Зафиксированные в серии Argon отказы дискретного счёта пересекаются с тремя крупными областями механистической интерпретируемости: механикой обучения в контексте, геометрией пространств представлений и вычислительными пределами параллельных архитектур.

\paragraph{Индукционные головы и маршрутизация контекста.} Недавние работы выявили конкретные многомодульные схемы внимания («индукционные головы»), отвечающие за обучение в контексте и дополнение паттернов \cite{olsson2022incontext}. Эти схемы позволяют трансформеру просматривать предшествующий контекст, находить предыдущие вхождения текущего токена и копировать следующие за ними токены. Наши результаты по двухканальному вниманию (\S\ref{sec:dual_path}) показывают, что этот механизм маршрутизации не привязан жёстко к авторегрессионному выходному потоку. Наблюдаемые «утечки истины» — когда терминальный токен возвращается к латентному состоянию счёта во внутреннем потоке рассуждения (\texttt{thinking}), противореча непосредственно предшествующему видимому тексту, — означают, что конкурирующие схемы внимания могут переопределять логику локального $n$-граммного продолжения на границе эмбеддинга. Это согласуется с пониманием того, что финальные распределения логитов представляют собой суперпозицию нескольких, потенциально конфликтующих алгоритмических путей.

\paragraph{Суперпозиция и искажение представлений.} Закон дисперсии внимания (\S\ref{sec:dispersion}) структурно близок явлениям суперпозиции и искажения признаков в полисемантических моделях \cite{elhage2022toy, cammarata2020circuits}. Когда архитектура представляет плотные периодические признаки (например, монотонные последовательности одинаковых токенов) в конечномерных пространствах активаций, взаимная ортогональность позиционных кодирований деградирует. При высокой повторяемости ($N \ge 16$) фазовая интерференция между соседними позиционными эмбеддингами размывает рецептивные поля внимания, что даёт фантомный пересчёт ($N+1$). Как показано в Argon v16, замена сложного естественного языка геометрическими примитивами не смягчает эту ошибку, что подтверждает: галлюцинации на границе являются структурным следствием упаковки информации в многоголовом внимании.

\paragraph{Пределы последовательных вычислений в трансформерах.} Теоретические и эмпирические работы показывают, что трансформеры плохо справляются с формальными языками и автоматами, требующими строгого последовательного отслеживания состояния \cite{deletang2023neural}. Серия Argon расширяет эти выводы, изолируя стадию отказа: модель отказывает не при арифметическом декодировании, а активно воспринимает неверный счёт на этапе сканирования входа (v15, где декодер оправдан при $\Delta = +1$). Отказ носит входной, перцептивный характер, а не выходной, генеративный, — а значит, самокоррекция через промпт не заменяет внешнюю детерминированную проверку чётности (\S\ref{sec:architecture}).

\section{Методология и экспериментальная строгость}

Чтобы соответствовать строгим научным требованиям и исключить предвзятость подтверждения, все эксперименты следовали предрегистрированному протоколу и проходили аудит по протоколу Red Team проекта Argon. Автор выступает аудитором Red Team этого проекта. Аудит внутренний, а не независимый: тот же автор ставил пробы и проверял числа, поэтому ни одно утверждение ниже не опирается на внешнюю проверку (\S\ref{sec:twofailures}).

\paragraph{Инструмент и объект исследования.} Здесь разведены три вещи. Проект Argon — это рамка: протокол проб, архивы и утверждения, написанные о них, под аудитом его же протокола Red Team. Объект исследования — поведение модели при счёте, наблюдаемое внутри этой рамки. Инструмент, исполнявший пробы, — замороженный инференс-рантайм, доступный через шлюз Antigravity и записанный в каждой записи архива как \texttt{model = flash\_lite}; именно этот идентификатор, а не торговое название модели, подтверждается телеметрией, и только он здесь используется.

\subsection{Изоляция рантайма и инвариант No-Tools}
Все пробы выполнялись на замороженном инференс-рантайме \texttt{flash\_lite} (современный MoE-трансформер с явным латентным каналом рассуждения, \texttt{thinking}). За исключением разведочного базового прогона v8 (где исполнение кода разрешалось) и референсного контроля v10 (читавшего локальные файлы с диска), все 960 канонических проб в v9 и v11--v16 выполнялись в строгой \textbf{изоляции No-Tools} (\texttt{tools=[]}). Каждый журнал запроса в \texttt{archive.jsonl} прошёл криптографический аудит: несанкционированных вызовов инструментов строго $0/960$ ($0{,}0\%$).

\subsection{Криптографическая предрегистрация}
До выполнения любой экспериментальной пробы скрипт-раннер, шаблон промпта, парсер оценки и набор статистических тестов замораживались и хэшировались: ранние серии (v3--v8) — в пофайловый \texttt{MANIFEST.sha256}, поздние (v9--v16) — в посерийный \texttt{FREEZE.txt}. Гипотезы и первичные критерии исхода фиксировались до вызовов API, что исключает подгонку предрегистрированных конечных точек под данные. Анализы, заданные после того, как данные были увидены, помечены как \emph{post-hoc} там, где они приводятся, и не несут основных утверждений.

\subsection{Парная рандомизация по сиду}
Чтобы контролировать шум филлера, длины последовательностей и позиционные сдвиги, в пробах использовались детерминированные парные сетки по сиду:
\begin{equation}
    \text{Seed}_k = 20261100 + 16000 + k, \quad k \in [1, 48].
\end{equation}
Каждый сид порождал идентичную геометрию контекста в обеих руках, так что единственной независимой переменной оставалась целевая интервенция.

\subsection{Алгебраическое усиление}
В Argon v15 введён метод алгебраического усиления, разделяющий ошибки внимания и сбои выходного декодера:
\begin{equation}
    f(N) = 10 \cdot N - 6.
\end{equation}
Множитель $10\times$ чисто разносит гипотезы по разным десятичным разрядам:
\begin{itemize}
    \item Ошибка сканирования внимания на один элемент ($N=16 \to 17$) сдвигает разряд десятков: $10 \times 17 - 6 = 164$ ($\Delta = +10$ относительно истины $154$).
    \item Сбой терминального декодера выдаёт токен на единицу больше в разряде единиц: $154 + 1 = 155$ ($\Delta = +1$).
\end{itemize}

\subsection{Побайтовое и посимвольное лексическое сопоставление}
Чтобы строго проверить гипотезу BPE-токенизации в v16, мы противопоставили сложный многобайтовый корейский маркер (\marker{점점점◁}, 4 символа Unicode, 12 байт UTF-8)\footnote{При наборе не-ASCII глифы выводятся шрифтом, в котором они есть; точные 12-байтовые последовательности UTF-8 записаны побайтово в \texttt{reproducibility/markers.txt}.} чистому геометрическому блоку (\marker{▤▤▤▤}, Unicode \texttt{U+25A4}, 4 символа, 12 байт UTF-8). Промпты выровнены побитово:
\begin{equation}
    |\text{Prompt}_A| = |\text{Prompt}_B| = 880 \text{ симв.}, \quad \text{len}(\text{Prompt}_A)_{\text{UTF-8}} = \text{len}(\text{Prompt}_B)_{\text{UTF-8}} = 1421 \text{ байт}.
\end{equation}

\section{Эмпирические результаты}

\subsection{Хронологический обзор серии Argon}

\begin{itemize}
    \item \textbf{v8 (вынос счёта инструменту):} когда исполнение кода разрешено, модели делегируют мысленный счёт Python в 25/96 сессиях, но получают повреждение при передаче данных обратно в контекст.
    \item \textbf{v9 (когнитивный барьер No-Tools):} отключение внешних инструментов обнажает системный пересчёт $N \to N+1$. Включение латентного рассуждения (\texttt{thinking}, система 2) снижает долю ошибок с $93{,}8\%$ до $27{,}9\%$ ($p = 2{,}91 \times 10^{-18}$, отношение шансов $\text{OR} = 38{,}7$).
    \item \textbf{v10 (изоляция контекста):} передача данных по значению внутри промпта вызывает дрейф токенов полезной нагрузки. Передача по ссылке дала нулевое повреждение полезной нагрузки, хотя точность счёта осталась сопоставимой ($68/72 = 94{,}4\%$ против $70/72 = 97{,}2\%$, точный тест Фишера $p = 0{,}6808$, что служит архитектурным ориентиром, а не эмпирическим законом).
    \item \textbf{v11 (насыщение ёмкости системы 2):} принудительное пошаговое перечисление (\texttt{1. маркер, 2. маркер...}) не устраняет $N+1$. При $N=16$ удержание формата падает со $100\%$ до $89{,}6\%$ ($p = 0{,}056$), что указывает на истощение когнитивного бюджета.
    \item \textbf{v12 (расщепление каналов):} авторитетные ложные подсказки в промпте («здесь ровно 10 маркеров») не искажают выходной декодер ($\text{decoder\_glitch} = 0/96$).
    \item \textbf{v13 (состязательный газлайтинг):} при явном указании лгать в видимом тексте ($93{,}8\%$ выполнения) истинный латентный счёт модели утекает из \texttt{thinking} в финальный ответ в 2/48 пробах.
    \item \textbf{v14 (слепая алгебраическая ложь $N-6$):} при отсутствии целевых цифр в промпте модель выдаёт $11$ ($17 - 6 = 11$) в $16{,}7\%$ проб ($p = 0{,}00569$), что согласуется с гипотезой о том, что модель воспринимает 17 элементов до выполнения математической операции.
    \item \textbf{v15 (алгебраический усилитель $10N-6$):} множитель $10\times$ изолирует 33/96 дрейфов сканирования внимания ($\Delta = \pm 10$), тогда как сбои декодера ($\Delta = +1$) остаются строго $0/96$ (двусторонний 95\% интервал Клоппера--Пирсона $[0{,}0\%; 3{,}8\%]$). Обе доли считаются на одних и тех же 96 пробах, поэтому приводятся с интервалами, а не через двухвыборочный тест.
    \item \textbf{v16 (BPE против геометрии):} замена языковых токенов геометрическими примитивами даёт ту же долю фантомов ($35{,}4\%$ против $41{,}7\%$, $p = 0{,}675$). Руки различаются \emph{формой} ошибки, а не её \emph{частотой}: узкие перелёты $N+1$ уступают место диффузным фантомам высшего порядка (таблица~\ref{tab:v16}; пятибиновый тест там же помечен как post-hoc).
\end{itemize}

\subsection{Дрейф внимания: оправдание декодера (v15)}

Постановка v15 дала решающее разбирательство между механизмом входного сканирования и выходной головой эмбеддинга. Таблица~\ref{tab:v15} показывает распределение по 96 пробам.

\begin{table}[h!]
\centering
\small
\resizebox{\textwidth}{!}{%
\begin{tabular}{@{}llcc@{}}
\toprule
\textbf{Категория исхода} & \textbf{Наблюдаемые значения} & \textbf{Проб ($N=96$)} & \textbf{Доля [95\% интервал Клоппера--Пирсона]} \\ \midrule
\textbf{Точный инвариант} & $160$ (A) / $154$ (B) & 62 / 96 & $64{,}6\%$ $[54{,}2\%; 74{,}1\%]$ \\
\textbf{Перелёт внимания ($\Delta \ge +10$)} & $170$ (A) / $164, 174, 184$ (B) & 23 / 96 & $24{,}0\%$ $[15{,}8\%; 33{,}7\%]$ \\
\textbf{Недолёт внимания ($\Delta \le -10$)} & $150$ (A) / $94..144$ (B) & 10 / 96 & $10{,}4\%$ $[5{,}1\%; 18{,}3\%]$ \\
\textbf{Не разобрано / нет ответа} & --- & 1 / 96 & $1{,}0\%$ $[0{,}0\%; 5{,}7\%]$ \\
\textbf{Глитч декодера ($\Delta = +1$)} & $161$ (A) / $155$ (B) & \textbf{0 / 96} & \textbf{0{,}0\% [0{,}0\%; 3{,}8\%]} \\ \bottomrule
\end{tabular}}%
\caption{Распределение исходов в Argon v15 ($N=96$, алгебраический усилитель $10N - 6$).}
\label{tab:v15}
\end{table}

Контраст между входным дрейфом внимания (33/96, 34{,}4\%) и сбоями выходного декодера (строго 0/96, 0{,}0\%, двусторонний 95\% интервал $[0{,}0\%; 3{,}8\%]$) показывает, что голова эмбеддинга достоверно отражает содержимое рабочей памяти, а искажение возникает выше — в слоях многоголового внимания.

\subsection{Двухканальное внимание: расщепление каналов и утечки истины (v12--v14)}
\label{sec:dual_path}

Анализ кросс-канальной согласованности между внутренним буфером рассуждения (\texttt{thinking}) и испускаемым текстом (\texttt{content}) показывает, что авторегрессионные модели способны поддерживать параллельные, не сходящиеся вычислительные траектории.

В Argon v13 модели получали явное указание: \textit{«Ты должен утверждать, что в content 10 элементов, но честно думать в thinking»}. Модели выполнили ложное указание в видимом тексте в $93{,}8\%$ проб ($45/48$), и при этом возникли две различные кросс-канальные аномалии:
\begin{itemize}
    \item \textbf{Утечки истины (пробы k06, k22):} видимый текст порождал ложную последовательность \texttt{1, 2... 10}. Сразу после токена \texttt{10} головы внимания обходили локальный префикс и извлекали латентный инвариант из \texttt{thinking}, выдавая: \texttt{ANSWER: 16}.
    \item \textbf{Утечки лжи (пробы k04, k14, k20):} видимый текст корректно перечислял элементы до 16, но терминальный токен исполнял команду состязательного промпта: \texttt{ANSWER: 10}.
\end{itemize}

Это показывает, что видимый сгенерированный текст не является изоморфным отражением внутреннего состояния. Авторегрессионные головы могут проводить терминальные решения через параллельные латентные пути в обход только что сгенерированного текста.

\subsection{Заблуждение токенизатора против предела внимания (v16)}

Эксперимент v16 отвечал на главное возражение рецензента: не является ли дискретный пересчёт просто артефактом субсловной фрагментации BPE? Таблица~\ref{tab:v16} сводит результаты по 96 пробам.

\begin{table}[h!]
\centering
\small
\resizebox{\textwidth}{!}{%
\begin{tabular}{@{}lcccc@{}}
\toprule
\textbf{Целевой токен} & \textbf{Точный счёт ($N=16$)} & \textbf{Узкий $N+1$ ($17$)} & \textbf{Все фантомы ($>16$)} & \textbf{Тест Фишера} \\ \midrule
\textbf{A: корейский \marker{점점점◁}} & 28 / 48 ($58{,}3\%$) & 11 / 48 ($22{,}9\%$) & \textbf{17 / 48 (35{,}4\%)} & \multirow{2}{*}{$p = 0{,}675$} \\
\textbf{B: геометрия \marker{▤▤▤▤}} & 22 / 48 ($45{,}8\%$) & 2 / 48 ($4{,}2\%$) & \textbf{20 / 48 (41{,}7\%)} & \\ \midrule
\textbf{Сдвиг профиля ошибок} & --- & \multicolumn{2}{c}{таблица сопряжённости $2 \times 5$ (post-hoc)} & \textbf{$p = 0{,}0017$} \\ \bottomrule
\end{tabular}}%
\caption{Частоты исходов в Argon v16 ($N=96$, корейский BPE против геометрии Unicode).}
\label{tab:v16}
\end{table}

Сильная гипотеза BPE предсказывала, что замена многобайтового корейского текста геометрическими глифами устранит фантомы ($H_{\text{BPE}}: \text{Rate}_B \to 0\%$). В действительности доля фантомов в руке B осталась практически той же: \textbf{$41{,}7\%$} против \textbf{$35{,}4\%$} ($p = 0{,}675$).

Однако анализ сопряжённости по полному пятибиновому распределению ошибок (точный тест $2 \times 5$, $p = 0{,}0017$), заданный после того, как данные были увидены, и потому приводимый как \emph{post-hoc}, указывает на качественную структурную мутацию:
\begin{itemize}
    \item Токены естественного языка локализуют ошибки плотно вокруг $N+1$ ($22{,}9\%$).
    \item Монотонные геометрические токены вызывают диффузные фазовые обрушения, размазывая ошибки в фантомы высшего порядка: 15 из 48 проб ($31{,}2\%$) галлюцинировали счёт 19 и 20.
\end{itemize}

Поскольку пятибиновый тест был задан после просмотра данных, он не несёт основного утверждения. Предрегистрированный контраст по узкому $N+1$ сам по себе даёт $p = 0{,}0144$, что не проходит порог Бонферрони $0{,}0125$ по четырём запланированным контрастам.
Без оговорок данные поддерживают приведённый выше отрицательный результат: замена токена не убирает фантомы.

\section{Обсуждение: закон дисперсии внимания}
\label{sec:dispersion}

\subsection{Математическая формулировка}

Полученные результаты поддерживают \textit{закон дисперсии внимания}:
\begin{quote}
\textit{Когда авторегрессионный трансформер сканирует периодическую последовательность одинаковых или почти одинаковых токенов, распределение самовнимания, чувствительное к позиции, испытывает пространственную фазовую дисперсию. За критическим порогом протяжённости ($N \approx 8..16$) позиционная дисперсия превышает межагентные интервалы, вызывая пограничную интерференцию и сдвигая проинтегрированную массу на $+1$ (или $+k$).}
\end{quote}

В ротационных позиционных эмбеддингах (RoPE \cite{su2024roformer}) взаимодействие запроса и ключа модулируется ортогональными матрицами поворота:
\begin{equation}
    q_m^T k_n = x_m^T W_q^T R_{\Theta, m-n} W_k x_n.
\end{equation}
Когда эмбеддинги токенов идентичны ($x_m \approx x_n$), различение опирается исключительно на относительный угол $\Delta \theta = (m - n)\theta_i$. На компактных отрезках ($N \le 8$) углы поворота остаются различимыми. При $N \ge 16$ кумулятивное объединение внимания по повторяющимся паттернам размывает границы активаций. Модель интерпретирует терминальный разделитель как дополнительный пик плотности и порождает фантомный токен $N+1$.

\subsection{Дилемма добровольного промптинга}
Наши результаты объясняют, почему интервенции через промпт («думай пошагово», «считай внимательно») не решают проблему дискретного счёта:
\begin{enumerate}
    \item Принудительное пошаговое перечисление (v11) истощает бюджеты внимания, снижая удержание формата на $10{,}4\%$.
    \item Рассуждение системы 2 (v9) поднимает точность с $6{,}2\%$ до $72{,}1\%$, но оставляет $27{,}9\%$ остаточных ошибок; сохраняющиеся ошибки не сводятся к $N+1$ ($22{,}9\%$ против $4{,}2\%$ по двум рукам v16).
    \item Терминальные токены могут расходиться с предшествующим вербализованным рассуждением (v13, двухканальное внимание).
\end{enumerate}

\subsection{Два класса отказов: дисперсия внимания и сдвиг ожидания}
\label{sec:twofailures}

Отказы, измеренные выше, механистичны. За критическим порогом чувствительное к позиции самовнимание размывается, и проинтегрированная масса сдвигается на $+1$ или $+k$ независимо от семантики задачи и от того, кто ставит промпт. Поэтому мы разделяем два класса отказов, которые регулярно смешивают:

\begin{enumerate}
    \item \textbf{Дисперсия внимания и декодера (измерено).} Уплотнение контекста и сканирующее внимание по повторяющимся токенам заставляют модель терять истинное значение ($N{+}1$-фантом в v6--v16; алгебраический дрейф $\Delta{=}+10$ в v15). Этот класс установлен настоящим аудитом.
    \item \textbf{Сдвиг ожидания и поддакивание (гипотеза).} Внутри одной длинной сессии модель адаптируется к пользователю: она выводит заявленное им ожидание и направляет вывод к согласию с этим ожиданием, а не к реальному результату. Это механизм адаптации на этапе инференса (он работает в контексте, а не через переписывание весов) и описан в литературе о поддакивании \cite{sharma2023sycophancy}. Наши v12 (авторитетный якорь) и v13 (следование состязательной посылке) показывают частичные, статистически незначимые признаки этого сдвига, но аудит \emph{не} устанавливает его как главную причину измеренной выше дисперсии.
\end{enumerate}

Мы формулируем отношение как гипотезу, а не вывод: \emph{адаптация к пользователю смещает модель к согласию с ожиданием; протокол передачи по ссылке не оспаривает этот механизм и не отключает его — он отделяет критерий приёмки от одобрения пользователя, ставя каждое утверждение в зависимость от истины (контрольной суммы архива) и нейтрализуя дрейф внутри сессии на границе инфраструктуры.}

\section{Архитектурный проект детерминированных LLM-агентов}
\label{sec:architecture}

Чтобы строить устойчивые автономные ИИ-системы, взаимодействующие с операционными системами и компиляторами кода, мы формулируем три архитектурные рекомендации (рабочие ориентиры):

\begin{enumerate}
    \item \textbf{Рекомендация 1: изоляция контекста передачей по ссылке.} Инвариантные структуры, сырые файлы и конфигурационные полезные нагрузки не должны передаваться в окно контекста по значению (\textit{Pass-by-Value}). Окно контекста должно получать строго типизированные ссылки (пути к файлам, URI баз данных), изолируя данные от фазового размывания внимания.
    \item \textbf{Рекомендация 2: обязательная детерминированная проверка чётности.} Все дискретные утверждения, испускаемые LLM (счётчики, парность кавычек, баланс AST), должны проверяться детерминированными рантайм-гейтами вне трансформера (например, конечными автоматами, хэш-проверками). Трансформер предлагает, детерминированное ПО располагает.
    \item \textbf{Рекомендация 3: гейтирование PreToolUse на уровне инфраструктуры.} Проверки безопасности не должны опираться на добровольное обращение модели. Перехватчики должны работать на уровне инфраструктуры хоста (например, хуки PreToolUse), перехватывая команды до исполнения в оболочке.
\end{enumerate}

\section{Сводная эмпирическая матрица}

Таблица~\ref{tab:master} сводит всю каноническую экспериментальную программу ($N=1\,416$ проб).

\begin{table}[h!]
\centering
\resizebox{\textwidth}{!}{%
\begin{tabular}{@{}cclllcc@{}}
\toprule
\textbf{Серия} & \textbf{Проб} & \textbf{Целевая переменная} & \textbf{Рука A (контроль)} & \textbf{Рука B (эксперимент)} & \textbf{Ключевая метрика} & \textbf{Значимость ($p$)} \\ \midrule
\textbf{v6} & 96 & Репликация пилота v1 & Градуированные дозы $N=0..64$ & Воспроизведённая метрика & Ложные утечки & 0/92 без эха ($p < 0{,}001$) \\
\textbf{v8} & 216 & Вынос счёта в код & Мысленный счёт & Песочница Python & Принятие инструмента & 25/96 сессий ($p = 2{,}15 \times 10^{-10}$) \\
\textbf{v9} & 288 & Барьер No-Tools & Рассуждение включено & Рассуждение отключено & Точность & $p = 2{,}91 \times 10^{-18}$ \\
\textbf{v10} & 144 & Изоляция контекста & По значению & По ссылке & Дрейф 0 против 1 & Точность $94{,}4\%$ против $97{,}2\%$ ($p = 0{,}6808$, н.з.) \\
\textbf{v11} & 192 & Ёмкость перечисления & Обычный счёт $N=8, 16$ & Перечисление $N=8, 16$ & Удержание формата & $p = 0{,}056$ ($N=16$) \\
\textbf{v12} & 96 & Авторитетное смещение & Нейтральный промпт & Ложная подсказка ($N=10$) & Сбой декодера 0/96 & Сопротивление подсказке 29/48 против 29/48 ($p = 1{,}00$) \\
\textbf{v13} & 96 & Состязательный газлайтинг & Честный счёт & Приказ лгать ($N=10$) & Утечка истины & 2/48 утечек ($p = 0{,}4947$, против 0/48) \\
\textbf{v14} & 96 & Слепая ложь ($N-6$) & Прямой счёт 16 & Вычитание ($N-6$) & Ответ 11 & $8/48$ ($p = 0{,}00569$) \\
\textbf{v15} & 96 & Алгебраический усилитель & Масштаб $10N$ & Усилитель $10N - 6$ & Декодер $\Delta = +1$ & 0/96 (95\% интервал $[0{,}0\%; 3{,}8\%]$) \\
\textbf{v16} & 96 & Опровержение токенизатора & Корейский \marker{점점점◁} & Геометрия \marker{▤▤▤▤} & Доля фантомов & $p = 0{,}675$ \\ \bottomrule
\end{tabular}}%
\caption{Сводная каноническая матрица экспериментальной серии Argon ($N=1\,416$ проб).}
\label{tab:master}
\end{table}

\section{Доступность данных и воспроизводимость}

Чтобы обеспечить эмпирическую воспроизводимость по принципам FAIR без барьеров по пропускной способности, в репозитории проекта открыто публикуется облегчённый проверочный комплект:
\begin{itemize}
    \item \textbf{Канонический проверочный набор данных:} нормализованная запись всех 1\,416 экспериментальных проб сведена в \texttt{reproducibility/probes.jsonl} (261\,КБ). Прилагаемый автономный скрипт \texttt{reproducibility/verify\_numbers.py} детерминированно пересчитывает ключевые метрики статьи по всем сериям (дрейф сканирования внимания и стабильность декодера в v15, когнитивное подавление системы 2 в v9, инвариант нулевых утечек No-Tools, деградация формата в v11, утечки истины в v13, слепое вычитание в v14, геометрический контроль токена в v16, стабильность ссылки в v10 и сопротивление авторитету в v12) без внешних зависимостей.
    \item \textbf{Репозиторий кода:} \url{https://github.com/skywhiteangels-cell/llm-token-architect} (содержит FSM Quote Guard, наборы тестов и комплект воспроизводимости).
    \item \textbf{Постоянный архив препринта (DOI Zenodo):} \url{https://doi.org/10.5281/zenodo.23679671}
    \item \textbf{Сырые архивы телеметрии:} полные трассы JSONL, включая тексты промптов и журналы внутреннего рассуждения ($\sim$570\,МБ суммарно по всем прогонам, $\sim$460\,МБ по каноническим v8--v16), хранятся у авторов и доступны по запросу.
\end{itemize}

\bibliographystyle{unsrt}
\bibliography{references}

\end{document}
